\subsection{\texorpdfstring{\(\mathfrak{h}\) 在 \(\mathfrak{g}\) 的包络代数中的交换子}{\textbackslash mathfrak\{h\} 在 \textbackslash mathfrak\{g\} 的包络代数中的交换子}}

设 U 是 g 的包络代数，\(V \subset U\) 是 h 的包络代数。代数 V 可以与 h 的对称代数 \(\mathbf{S}(\mathfrak{h})\) 等同，也可以与 \(h^{*}\) 上的多项式函数代数等同。用 \(\alpha_{1},\ldots,\alpha_{n}\) 记两两不同的正根。设 \((H_{1},\ldots,H_{l})\) 是 h 的一个基。由 Poincaré-Birkhoff-Witt 定理，诸元素

\[
u ((q _ {i}), (m _ {i}), (p _ {i})) = X _ {- \alpha_ {1}} ^ {q _ {1}} \dots X _ {- \alpha_ {n}} ^ {q _ {n}} H _ {1} ^ {m _ {1}} \dots H _ {l} ^ {m _ {l}} X _ {\alpha_ {1}} ^ {p _ {1}} \dots X _ {\alpha_ {n}} ^ {p _ {n}}
\]

\((q_{i},m_{i},p_{i})\) 为整数 \(\geq 0\) ）构成向量空间 U 的一个基。对一切 \(h\in \mathfrak{h}\) ，我们有

\[
[ h, u ((q _ {i}), (m _ {i}), (p _ {i})) ] = ((p _ {1} - q _ {1}) \alpha_ {1} + \dots + (p _ {n} - q _ {n}) \alpha_ {n}) (h) u ((q _ {i}), (m _ {i}), (p _ {i})).\tag{3}
\]

向量空间 U 在伴随表示下是一个 g-模（从而也是一个 h-模）。若 \(\lambda \in h^{*}\) ，则子空间 \(U^{\lambda}\) 与 \(U_{\lambda}\) 已有定义（第七章 §1 第 3 节）；公式 (3) 表明 \(U^{\lambda} = U_{\lambda}\) ，且 \(U = \bigoplus_{\lambda \in Q} U^{\lambda}\) （其中 Q 是 R 的根权格）。特别地，\(U^{0}\) 是 h 或 V 在 U 中的交换子。

引理 3. 置 \(\mathrm{L} = (\mathfrak{n}_{-}\mathrm{U})\cap \mathrm{U}^{0}\)

\begin{enumerate}
\def\labelenumi{(\roman{enumi})}
\item
  我们有 \(\mathrm{L} = (\mathrm{U}\mathfrak{n}_{+})\cap \mathrm{U}^{0}\) ，且 \(\mathrm{L}\) 是 \(\mathrm{U}^0\) 的一个双边理想。
\item
  我们有 \(U^{0} = V \oplus L\) 。
\end{enumerate}

显然 \(\mathfrak{n}_{-}\mathrm{U}\) （相应地 \(\mathrm{U}\mathfrak{n}_{+}\) ）是这样的元素 \(u((q_{i}),(m_{i}),(p_{i}))\) 的线性组合组成的集合：它们使 \(\sum q_{i}>0\) （相应地 \(\sum p_{i}>0\) ）。另一方面

\[
u \left(\left(q _ {i}\right), \left(m _ {i}\right), \left(p _ {i}\right)\right) \in \mathrm{U} ^ {0} \Longleftrightarrow p _ {1} \alpha_ {1} + \dots + p _ {n} \alpha_ {n} = q _ {1} \alpha_ {1} + \dots + q _ {n} \alpha_ {n}.
\]

这蕴含 \((\mathfrak{n}_{-}\mathrm{U})\cap\mathrm{U}^{0}=(\mathrm{U}\mathfrak{n}_{+})\cap\mathrm{U}^{0}\) 。最后，\((\mathfrak{n}_{-}\mathrm{U})\cap\mathrm{U}^{0}\) （相应地 \((\mathrm{U}\mathfrak{n}_{+})\cap\mathrm{U}^{0}\) ）是 \(U^{0}\) 的一个右（相应地左）理想，由此得 (i)。此外，元素 \(u((q_{i}),(m_{i}),(p_{i}))\) 若属于 \(U^{0}\) ，则它属于 V（相应地属于 L）当且仅当 \(p_{1}=\cdots=p_{n}=q_{1}=\cdots=q_{n}=0\) （相应地 \(p_{1}+\cdots+p_{n}+q_{1}+\cdots+q_{n}>0\) ），由此得 (ii)。证毕。

鉴于引理 3，\(U^{0}\) 到 V 上的、以 L 为核的投影是一个代数同态。它称为从 \(U^{0}\) 到 V 的 Harish-Chandra 同态（相对于 B 的）。回忆 V 可以与 \(h^{*}\) 上的多项式函数代数等同。

命题 7. 设 \(\lambda \in \mathfrak{h}^*\) ，E 是一个 \(\mathfrak{g}\) -模，由权为 \(\lambda\) 的一个本原元生成，\(\chi\) 是 E 的中心特征标，\(\varphi\) 是从 \(\mathrm{U}^0\) 到 V 的 Harish-Chandra 同态。则 \(\chi(z) = (\varphi(z))(\lambda)\) 对 U 的中心中的所有 \(z\) 成立。

设 v 是 E 的权为 \(\lambda\) 的本原元，z 是 U 的中心的一个元素。存在 \(u_{1},\ldots,u_{p}\in U\) 与 \(n_{1},\ldots,n_{p}\in\mathfrak{n}_{+}\) 使 \(z=\varphi(z)+u_{1}n_{1}+\cdots+u_{p}n_{p}\) 成立。则

\[
\chi (z) v = z v = \varphi (z) v + u _ {1} n _ {1} v + \dots + u _ {p} n _ {p} v = \varphi (z) v = (\varphi (z)) (\lambda) v.
\]

推论. 设 \(\langle\cdot,\cdot\rangle\) 是 g 上的一个非退化不变对称双线性型，C 是与 \(\langle\cdot,\cdot\rangle\) 相伴的 Casimir 元。也用 \(\langle\cdot,\cdot\rangle\) 记 \(h^{*}\) 上的一个逆型，即 \(\langle\cdot,\cdot\rangle\) 在 h 上的限制的逆型（§2 第 3 节 命题 5）。则 \(\chi(\mathrm{C})=\langle\lambda,\lambda+2\rho\rangle\) ，其中 \(\rho=\frac{1}{2}\sum_{\alpha\in R_{+}}\alpha\) 。

我们回忆 §2 第 3 节 命题 6 的记号。我们有

\[
\begin{array}{r} \mathrm{C} = \sum_ {\alpha \in \mathrm{R} _ {-}} \frac {1}{\langle X _ {\alpha} , X _ {- \alpha} \rangle} X _ {\alpha} X _ {- \alpha} + \sum_ {\alpha \in \mathrm{R} _ {+}} \frac {1}{\langle X _ {\alpha} , X _ {- \alpha} \rangle} X _ {- \alpha} X _ {\alpha} \\ + \sum_ {\alpha \in \mathrm{R} _ {+}} \frac {1}{\langle X _ {\alpha} , X _ {- \alpha} \rangle} [ X _ {\alpha}, X _ {- \alpha} ] + \sum_ {i \in I} H _ {i} H _ {i} ^ {\prime} \end{array}
\]

故

\[
\varphi (\mathrm{C}) = \sum_ {\alpha \in \mathrm{R} _ {+}} \frac {1}{\langle X _ {\alpha} , X _ {- \alpha} \rangle} [ X _ {\alpha}, X _ {- \alpha} ] + \sum_ {i \in I} H _ {i} H _ {i} ^ {\prime}.
\]

由命题 7，

\[
\chi (\mathrm{C}) = \sum_ {\alpha \in \mathrm{R} _ {+}} \frac {1}{\langle X _ {\alpha} , X _ {- \alpha} \rangle} \lambda ([ X _ {\alpha}, X _ {- \alpha} ]) + \sum_ {i \in I} \lambda (H _ {i}) \lambda (H _ {i} ^ {\prime}).
\]

设 \(h_{\lambda}\) 是 \(\mathfrak{h}\) 中使 \(\langle h_{\lambda}, h \rangle = \lambda(h)\) 对一切 \(h \in \mathfrak{h}\) 成立的元素。由 §2 第 2 节 命题 1，

\[
\lambda \left(\frac {1}{\langle X _ {\alpha} , X _ {- \alpha} \rangle} [ X _ {\alpha}, X _ {- \alpha} ]\right) = \left\langle h _ {\lambda}, \frac {1}{\langle X _ {\alpha} , X _ {- \alpha} \rangle} [ X _ {\alpha}, X _ {- \alpha} ] \right\rangle = \alpha (h _ {\lambda}) = \langle \lambda , \alpha \rangle .
\]

于是

\[
\chi (\mathrm{C}) = \left(\sum_ {\alpha \in \mathrm{R} _ {+}} \langle \lambda , \alpha \rangle\right) + \langle \lambda , \lambda \rangle = \langle \lambda , \lambda + 2 \rho \rangle .
\]

